\documentclass[11pt]{article}
\usepackage[a4paper,margin=29mm]{geometry}
\usepackage{amsmath,amssymb,amsthm,mathtools}
\usepackage[T1]{fontenc}
\usepackage{lmodern}
\usepackage[colorlinks=true,linkcolor=blue,citecolor=blue,urlcolor=blue]{hyperref}
\newtheorem{theorem}{Theorem}
\newtheorem{lemma}[theorem]{Lemma}
\newtheorem{proposition}[theorem]{Proposition}
\newtheorem{remark}[theorem]{Remark}
\newcommand{\R}{\mathbb R}
\newcommand{\C}{\mathbb C}
\newcommand{\diag}{\operatorname{diag}}
\newcommand{\Id}{I}
\newcommand{\norm}[1]{\left\lVert#1\right\rVert}
\newcommand{\abs}[1]{\left\lvert#1\right\rvert}
\DeclareMathOperator{\Imat}{Im}
\DeclareMathOperator{\sigmamin}{\sigma_{min}}
\title{Logarithmic convexity for a coupled system\\with two subdiffusion orders}
\author{Kaiwen Jin\\\small\texttt{yui.ui.siki@gmail.com}}
\date{October 15, 2026}
\begin{document}
\maketitle
\begin{abstract}
We prove a logarithmic convexity estimate for a pair of coupled Caputo
subdiffusion equations with different time orders. For every bounded smooth
domain, every $0<\alpha_1<\alpha_2<1$, and every final time $T>0$, the product
$L^2$ norm at time $t$ is bounded by a constant times the geometric
interpolation of the initial and final norms, with exponent $t/T$.
The estimate holds for arbitrary $L^2$ initial pairs and includes domains
with growing low-frequency modes. This gives an affirmative answer to the
Dirichlet norm-interpolation question for the system posed by Chorfi.
The proof uses a positive matrix integral representation to establish
finite-time invertibility and uniform bounds for all large spatial
frequencies.
\end{abstract}

\section{The problem and the main result}\label{sec:main}
Subdiffusion equations describe evolution with a memory term in place of an
ordinary time derivative. Their backward problem asks whether an earlier
state can be recovered from a later observation. A logarithmic convexity
estimate bounds an intermediate state by a geometric interpolation of the
initial and final norms; in particular, it implies backward uniqueness.
For a single Caputo order, Chorfi, Maniar and Yamamoto
\cite[Theorem 1.1]{CMY} proved such an estimate for self-adjoint evolution
equations. Coupled systems with a common time order admit a spectral treatment
of their forward and backward problems \cite{FLL}. Mixed-order coupled
systems also have a well-developed mild-solution theory \cite{LHL}, but that
existence theory does not itself give interpolation between the endpoint
norms.

The distinction between a common order and two different orders is essential
for the proof: a rotation that diagonalizes the reaction matrix mixes the two
time derivatives. Chorfi \cite[Section 5, problem 1]{Chorfi} asks whether
logarithmic convexity extends to the particular mixed-order system
\begin{equation}\label{eq:PDE}
 \begin{aligned}
 D_t^{\alpha_1}u_1&=\Delta u_1+u_1+u_2,\\
 D_t^{\alpha_2}u_2&=\Delta u_2+u_1+u_2.
 \end{aligned}
\end{equation}
We consider homogeneous Dirichlet conditions on a bounded smooth domain
$\Omega\subset\R^n$, with arbitrary initial pair
$U_0\in H:=L^2(\Omega;\C^2)$ and $U=(u_1,u_2)$.
The norm is
$\norm U^2=\norm{u_1}_{L^2(\Omega)}^2+\norm{u_2}_{L^2(\Omega)}^2$.
For an absolutely continuous scalar function, the Caputo convention is
\[
 D_t^\alpha h(t)=\frac1{\Gamma(1-\alpha)}
       \int_0^t(t-s)^{-\alpha}h'(s)\,ds.
\]
For $L^2$ initial data, we use the continuous $H$ representative of the
componentwise mild solution, equivalently the fractional diffusion Duhamel
formulation in \cite[Definition 1]{LHL}. The real case follows by restriction.

\begin{theorem}\label{thm:main}
Let $\Omega\subset\R^n$ be any bounded smooth domain,
$0<\alpha_1<\alpha_2<1$, and $T>0$.
There is a constant $C=C(\Omega,\alpha_1,\alpha_2,T)>0$ such that every
homogeneous Dirichlet mild solution of \eqref{eq:PDE} with
$U(0)=U_0\in H$ satisfies
\begin{equation}\label{eq:main}
 \norm{U(t)}\le C\norm{U_0}^{1-t/T}\norm{U(T)}^{t/T},
 \qquad 0<t<T.
\end{equation}
The constant is independent of the initial pair and of the intermediate
time. In particular, $U(T)=0$ implies $U_0=0$.
\end{theorem}
This is the product-$L^2$ Dirichlet formulation of the question, including all
initial pairs and the exact interpolation exponent $t/T$; it is also the
formulation of AIM Problem~402.\footnote{The frozen problem statement is
\href{https://github.com/MColbrook/AIM/blob/a153f469ee0ea64e876b2ad97c79d252c19a3e22/problems/402-mixed-order-fractional-logarithmic-convexity.md}{AIM Problem 402 at revision \texttt{a153f469ee0e}}.}
No sign condition or additional initial regularity is required.

We keep the Dirichlet spatial eigenbasis and treat each pair of coefficients
as a two-dimensional system. A positive matrix integral representation of its
resolvent proves that the terminal map is invertible, including the finitely
many modes that may grow. For large spatial eigenvalues, the same
representation gives an upper smoothing bound and a matching lower bound for
the smallest singular value at the final time. These bounds reduce
\eqref{eq:main} to scalar interpolation and H\"older's inequality.
The constant is not asserted to be sharp or equal to one. The argument uses
the common Laplacian and the symmetric constant coupling in \eqref{eq:PDE};
it does not treat general spatially varying coupling.

The modal representation and invertibility are proved in Section
\ref{sec:modal}. Section \ref{sec:uniform} establishes the uniform estimates
and proves the theorem. In the proof we write $a=\alpha_1$ and
$b=\alpha_2$. Matrix norms and vector lengths are Euclidean; inequalities
between real symmetric matrices refer to their quadratic forms.


\section{The modal resolvent and finite-time invertibility}\label{sec:modal}
Let $0<a<b<1$ be fixed. A Dirichlet eigenvalue $\lambda>0$ gives the
two-dimensional system
\begin{equation}\label{eq:mode}
 D_t^a y_1=-qy_1+y_2,\qquad
 D_t^b y_2=y_1-qy_2,\qquad q=\lambda-1>-1.
\end{equation}
Denote its fundamental matrix by $F_q(t)$, with $F_q(0)=\Id$, and put
\[
 A_q=\begin{pmatrix}q&-1\\-1&q\end{pmatrix},\qquad
 D(z)=\diag(z^a,z^b),\qquad R_q(z)=(D(z)+A_q)^{-1}.
\]
All complex powers use the branch $\arg z\in(-\pi,\pi)$.
The integral form of \eqref{eq:mode} is
$y_i(t)=y_i(0)-I^{\alpha_i}(A_qy)_i(t)$, where
$\alpha_1=a$, $\alpha_2=b$ and
\[
 I^\gamma f(t)=\frac1{\Gamma(\gamma)}
       \int_0^t(t-s)^{\gamma-1}f(s)\,ds.
\]
It has a unique continuous solution of exponential order. Indeed, in the
component maximum norm weighted by $e^{-\omega t}$, $\omega\ge1$, its
Volterra operator has norm at most $(|q|+1)\omega^{-a}$.
For large $\omega$, the Neumann series converges. Local uniqueness also
follows by iteration, since on $[0,T]$ the row kernels are bounded by
$C_Tt^{a-1}$ and their $m$th convolution integrals are bounded by
$(C_Tt^a)^m/\Gamma(1+ma)$. Taking Laplace transforms gives
\begin{equation}\label{eq:Fhat}
 \widehat F_q(z)=\frac1zR_q(z)D(z)
               =\frac1z(\Id-R_q(z)A_q)
\end{equation}
for positive real $z$ sufficiently large. In particular, the different
initial factors $z^{a-1}$ and $z^{b-1}$ have been retained.

The following positive matrix density will distinguish invertibility from
mere positivity of the entries of a solution.
For $r>0$, define
\begin{align*}
 M_q(r)&=A_q+\diag(r^ae^{i\pi a},r^be^{i\pi b}),\\
 B(r)&=\diag(r^a\sin\pi a,r^b\sin\pi b),\\
 W_q(r)&=-\frac1\pi\Imat R_q(-r+i0)
       =\frac1\pi M_q(r)^{-*}B(r)M_q(r)^{-1}.
\end{align*}
Here $X^*$ denotes conjugate transpose and $\Imat X=(X-X^*)/(2i)$.
The inverse identity $\Imat(M^{-1})=-M^{-*}(\Imat M)M^{-1}$ proves the
last equality. Since $B(r)>0$, the density $W_q(r)$ is strictly positive
definite. It is real symmetric: $R_q(z)$ is complex symmetric, so its
Hermitian imaginary part is its entrywise imaginary part.

\begin{lemma}\label{lem:representation}
For every $q>-1$, the modal resolvent has the representation
\begin{equation}\label{eq:cut}
 R_q(z)=\int_0^\infty\frac{W_q(r)}{z+r}\,dr
       +\begin{cases}K_q/(z-\rho_q),&-1<q<1,\\0,&q\ge1.
         \end{cases}
\end{equation}
In the first case, $\rho_q>0$ is the unique positive pole and $K_q$ is a
positive semidefinite matrix of rank one. The integral is absolutely
convergent away from the cut and the pole.
\end{lemma}

\begin{proof}
If $\Im z>0$, then $\Imat D(z)>0$. Thus $D(z)+A_q$ has no kernel;
the lower half-plane follows by conjugation. On the positive real axis
$D(s)+A_q$ increases strictly in the positive definite ordering. Its
eigenvalues start at $q-1,q+1$ and tend to positive infinity. It therefore
has one zero, and only one, when $-1<q<1$, and none when $q\ge1$.
At the zero $\rho_q$, let $v_q$ be a real unit kernel vector. The other
eigenvalue is positive, and
\[
 K_q=\mathop{\rm Res}_{z=\rho_q}R_q(z)
     =\frac{v_qv_q^T}{v_q^TD'(\rho_q)v_q}.
\]
The denominator is positive, so this is a simple pole with the stated
residue.

Apply the Cauchy formula entrywise to a keyhole contour about the negative
real axis, retaining this pole when present. For fixed $q$,
$R_q(z)=O(|z|^{-a})$ at infinity, uniformly in the argument, by factoring
out $D(z)$. It is bounded at zero if $q\ne1$, since $q=-1$ is excluded.
At $q=1$ its explicit inverse is
\begin{equation}\label{eq:critical-resolvent}
 R_1(z)=\frac1{z^{a+b}+z^a+z^b}
        \begin{pmatrix}z^b+1&1\\1&z^a+1\end{pmatrix}
        =O(|z|^{-a})\quad(z\to0).
\end{equation}
The small and large contour integrals consequently vanish. The jump
between the two banks of the cut is $-2\pi iW_q(r)$, giving
\eqref{eq:cut} with the displayed residue sign.
At infinity $W_q(r)=O(r^{-a})$; at zero it is $O(r^a)$ if $q\ne1$ and
$O(r^{-a})$ if $q=1$. These bounds also prove the asserted absolute
convergence. No bound uniform in $q$ is needed in this lemma.
\end{proof}

Define the time kernel and its integral by
\begin{equation}\label{eq:time-kernel}
 Q_q(t)=\int_0^\infty e^{-rt}W_q(r)\,dr
       +\begin{cases}e^{\rho_qt}K_q,&q<1,\\0,&q\ge1,
        \end{cases}
 \qquad G_q(t)=\int_0^tQ_q(s)\,ds.
\end{equation}
The kernel is strictly positive definite. Its singularity at zero is at
most $C_qt^{a-1}$, so it is locally integrable. Absolute Fubini and
\eqref{eq:cut} identify its Laplace transform with $R_q$ to the right of
the pole. Equation \eqref{eq:Fhat} and uniqueness of the Laplace transform
therefore yield
\begin{equation}\label{eq:F-time}
 F_q(t)=\Id-G_q(t)A_q.
\end{equation}

\begin{proposition}\label{prop:invertibility}
The matrix $F_q(t)$ is invertible for every $q>-1$ and every finite
$t\ge0$.
\end{proposition}

\begin{proof}
Suppose first that $q\ne1$. The endpoint bounds in Lemma
\ref{lem:representation} ensure convergence of
\[
 P_q(t)=\int_0^\infty e^{-rt}\frac{W_q(r)}r\,dr>0
 \qquad(t\ge0).
\]
For $q<1$, evaluating \eqref{eq:cut} at zero and integrating
\eqref{eq:time-kernel} give, respectively,
\[
 \begin{aligned}
 A_q^{-1}&=P_q(0)-K_q/\rho_q,\\
 G_q(t)&=P_q(0)-P_q(t)+(e^{\rho_qt}-1)K_q/\rho_q.
 \end{aligned}
\]
For $q>1$, the same identities hold with the pole terms omitted.
Their constant terms cancel in \eqref{eq:F-time}, giving the real symmetric
matrix
\begin{equation}\label{eq:H}
 H_q(t):=F_q(t)A_q^{-1}
 =P_q(t)-\begin{cases}e^{\rho_qt}K_q/\rho_q,&q<1,\\0,&q>1.
            \end{cases}
\end{equation}
For $q>1$, both $H_q(t)$ and $A_q$ are positive definite, and the result
follows. For $-1<q<1$, the matrix $A_q^{-1}$ has one eigenvalue of each
sign. Since $H_q(t)=A_q^{-1}-G_q(t)$, its negative direction remains
strictly negative. A nonzero vector in $\ker K_q$ has strictly positive
quadratic form under $H_q(t)$, by \eqref{eq:H}. Hence $H_q(t)$ has one
eigenvalue of each sign and is nonsingular.

For $q=1$, put $w=(1,-1)^T$, so $A_1=ww^T$. From
\eqref{eq:critical-resolvent},
\[
 w^TR_1(s)w=\frac{s^a+s^b}{s^{a+b}+s^a+s^b}\longrightarrow1
       \qquad(s\downarrow0).
\]
Writing $k(r)=w^TW_1(r)w>0$, the representation \eqref{eq:cut} and
monotone convergence imply $\int_0^\infty k(r)/r\,dr=1$.
The rank-one determinant formula in \eqref{eq:F-time} now gives
\[
 \det F_1(t)=1-w^TG_1(t)w
           =\int_0^\infty e^{-rt}\frac{k(r)}r\,dr>0.
\]
This also covers a domain having the critical Dirichlet eigenvalue
$\lambda=2$. At $t=0$, invertibility is immediate from $F_q(0)=\Id$.
\end{proof}

\section{Uniform estimates and the interpolation inequality}\label{sec:uniform}
Invertibility of each fixed mode alone is insufficient for the desired
estimate. We next control all large eigenvalues with the same constants.
Throughout this section, $T>0$ remains fixed.

\begin{lemma}\label{lem:uniform}
There are $M=M(a,b,T)>0$ and $d=d(a,b,T)>0$ such that, for every $q\ge2$,
\begin{equation}\label{eq:uniform}
 \norm{F_q(t)}\le\frac{M}{1+qt^b}\quad(0<t\le T),\qquad
 \abs{F_q(T)v}\ge\frac{d}{\sqrt3q}\abs v\quad(v\in\C^2).
\end{equation}
\end{lemma}

\begin{proof}
The eigenvalues of $A_q$ are $q-1,q+1$, so its condition number is at most
three. The positive integral in \eqref{eq:H} gives
\[
 S_q(t):=A_q^{1/2}H_q(t)A_q^{1/2}
       =A_q^{1/2}F_q(t)A_q^{-1/2}.
\]
This is positive definite, decreases in the ordering of symmetric matrices,
and satisfies $S_q(0)=\Id$. For $s>0$, set
$E=A_q^{-1/2}D(s)A_q^{-1/2}>0$. Equation \eqref{eq:Fhat} yields
\[
 \widehat S_q(s)=\frac1s\bigl(\Id-(\Id+E)^{-1}\bigr)
 =\frac1sE(\Id+E)^{-1}
 \le\frac{\max(s^a,s^b)}{(q-1)s}\Id.
\]
On the other hand, monotonicity gives
\[
 \widehat S_q(1/t)\ge\int_0^te^{-s/t}S_q(t)\,ds
                 =t(1-e^{-1})S_q(t).
\]
Since also $S_q(t)\le\Id$, returning through the similarity gives
\begin{equation}\label{eq:upper-intermediate}
 \norm{F_q(t)}\le\sqrt3\min\left(1,
   \frac{\max(t^{-a},t^{-b})}{(1-e^{-1})(q-1)}\right).
\end{equation}
Let $L_T=\max(1,T^{b-a})$. For $0<t\le T$ we have
$\max(t^{-a},t^{-b})\le L_Tt^{-b}$, and $q-1\ge q/2$.
Using $\min(1,c/x)\le(1+c)/(1+x)$ for $c,x>0$ in
\eqref{eq:upper-intermediate} proves the first estimate, for example with
$M=\sqrt3(1+2L_T/(1-e^{-1}))$.

For the lower estimate, recall the positive density
\[
 W_q(r)=\frac1\pi M_q(r)^{-*}B(r)M_q(r)^{-1}.
\]
Here the factors are
\[
 \begin{aligned}
 M_q(r)&=A_q+\operatorname{diag}(r^ae^{i\pi a},r^be^{i\pi b}),\\
 B(r)&=\operatorname{diag}(r^a\sin\pi a,r^b\sin\pi b).
 \end{aligned}
\]
Set $\eta=\min(\sin\pi a,\sin\pi b)>0$ and $L=1+2^b$. For $1\le r\le2$,
$B(r)\ge\eta\Id$ and $\norm{M_q(r)}\le q+L$. Therefore
\[
 W_q(r)\ge\frac{\eta}{\pi(q+L)^2}\Id.
\]
Restricting the positive integral for $H_q(T)$ to this interval gives
\[
 S_q(T)\ge\frac{\eta e^{-2T}(q-1)}{2\pi(q+L)^2}\Id
         \ge\frac dq\Id,\qquad
 d=\frac{\eta e^{-2T}}{4\pi(1+L/2)^2}.
\]
Here $q-1\ge q/2$ and $q+L\le q(1+L/2)$ were used.
Since $F_q(T)=A_q^{-1/2}S_q(T)A_q^{1/2}$, its smallest singular value is
at least $(d/q)\sqrt{(q-1)/(q+1)}\ge d/(\sqrt3q)$.
\end{proof}

We now deduce the modal interpolation estimate. Put $\theta=t/T$ and
\[
 h_q=(1+qT^b)^{-1},\qquad m=\min(1,dT^b/\sqrt3)>0.
\]
Lemma \ref{lem:uniform} implies $|F_q(T)v|\ge mh_q|v|$ for $q\ge2$.
The function $-\log(1+qt^b)$ is convex on $[0,T]$ and vanishes at zero.
Consequently,
\[
 (1+qt^b)^{-1}\le h_q^\theta,
\]
and the same lemma gives
\begin{equation}\label{eq:modal-interpolation}
 |F_q(t)v|\le \frac Mm|v|^{1-\theta}|F_q(T)v|^\theta
 \qquad(q\ge2,\ 0<t<T).
\end{equation}
The zero vector is included directly. No simultaneous diagonalization of
the two time derivatives is used.

\begin{proof}[Proof of Theorem \ref{thm:main}]
Let $(\lambda_j,\phi_j)$ be an orthonormal Dirichlet eigensystem of
$-\Delta$, and set $q_j=\lambda_j-1$ and
$v_j=((u_1(0),\phi_j),(u_2(0),\phi_j))^T$.
Projection of the componentwise mild equations gives
\eqref{eq:mode} in its fractional integral form. Thus its uniqueness gives
\begin{equation}\label{eq:series}
 U(t)=\sum_{j=1}^\infty\phi_jF_{q_j}(t)v_j.
\end{equation}
This uses only a bounded spectral coefficient functional on the mild
integrals. It does not apply $\Delta$ to $U_0$. Equivalently one can start
with the mild formulation in \cite[Definition 1 and Theorem 1]{LHL} and
project it. Lemma \ref{lem:uniform} and the finitely many lower modes give
a uniform operator bound on $[0,T]$; hence the series converges in product
$L^2$ and is continuous there, including at $t=0$, by dominated convergence.

There are only finitely many indices with $q_j<2$. If there are any, define
\[
 M_0=\max_{q_j<2}\max_{0\le s\le T}\norm{F_{q_j}(s)},\qquad
 m_0=\min_{q_j<2}\sigmamin(F_{q_j}(T)).
\]
Continuity and Proposition \ref{prop:invertibility} give $M_0<\infty$
and $m_0>0$. These modes satisfy \eqref{eq:modal-interpolation} with
constant $M_0/\min(1,m_0)$, since $m_0^\theta\ge\min(1,m_0)$.
Let $C$ be the larger of this constant and $M/m$, omitting the first
constant if the low-frequency set is empty. Then $C$ depends only on
$\Omega,a,b,T$, and Parseval and H\"older give
\begin{align*}
 \norm{U(t)}^2
 &\le C^2\sum_j |v_j|^{2(1-\theta)}|F_{q_j}(T)v_j|^{2\theta}\\
 &\le C^2\left(\sum_j|v_j|^2\right)^{1-\theta}
          \left(\sum_j|F_{q_j}(T)v_j|^2\right)^\theta\\
 &=C^2\norm{U_0}^{2(1-\theta)}\norm{U(T)}^{2\theta}.
\end{align*}
Taking square roots proves the result for every $U_0$ and $0<t<T$.
If $U(T)=0$, Proposition \ref{prop:invertibility} shows that every $v_j=0$,
so the same conclusion holds without a division by zero.
\end{proof}

\section*{AI assistance}
OpenAI Codex assisted with the literature search, development and checking
of the proof, and preparation of this manuscript. This assistance does not
replace independent mathematical review.

\begin{thebibliography}{9}
\bibitem{Chorfi}
S.-E. Chorfi, \emph{Logarithmic convexity of evolution equations and application
to inverse problems}, arXiv:2506.19954v1 (2025),
\href{https://arxiv.org/abs/2506.19954v1}{arXiv:2506.19954v1}.

\bibitem{CMY}
S. E. Chorfi, L. Maniar and M. Yamamoto, \emph{The backward problem for
time fractional evolution equations}, arXiv:2211.16493v1 (2022),
\href{https://arxiv.org/abs/2211.16493v1}{arXiv:2211.16493v1}.

\bibitem{LHL}
Z. Li, X. Huang and Y. Liu, \emph{Initial-boundary value problems for coupled
systems of time-fractional diffusion equations}, Fractional Calculus and
Applied Analysis \textbf{26} (2023), 533--566,
\href{https://doi.org/10.1007/s13540-023-00149-0}{doi:10.1007/s13540-023-00149-0}.

\bibitem{FLL}
D. Feng, Y. Liu and S. Lu, \emph{Forward and backward problems for coupled
subdiffusion systems}, arXiv:2407.00588v2 (2025),
\href{https://arxiv.org/abs/2407.00588v2}{arXiv:2407.00588v2}.
\end{thebibliography}
\end{document}
